\documentclass[12pt]{article}
\usepackage[utf8]{inputenc}
\usepackage{amsmath, amssymb}
\usepackage{xcolor}
\usepackage{geometry}
\usepackage{hyperref}
\usepackage{fancyhdr}
\usepackage{enumitem}
\usepackage{minted}
\usepackage{booktabs}
\usepackage{tikz}
\usetikzlibrary{shapes, arrows, positioning}

\geometry{margin=1in}
\hypersetup{colorlinks=true, linkcolor=blue, urlcolor=cyan}

\pagestyle{fancy}
\fancyhf{}
\fancyhead[L]{\textbf{\TOPICTITLE}}
\fancyhead[R]{\thepage}

% ------------------------------- 
% Topic Metadata
% ------------------------------- 
\newcommand{\TOPICTITLE}{Function Programming SKills Introduction}
\title{\TOPICTITLE\\\large Study-Ready Notes}
\author{Compiled by Andrew Photinakis}
\date{\today}

\setlength{\headheight}{15pt}

\begin{document}
\maketitle
\tableofcontents
\newpage

% \section{Function Programming SKills Introduction - Haskell}
\section{Understanding \texttt{foldr} in Haskell}

\subsection{The Abstract Pattern}
In Haskell, \texttt{foldr} (right fold) is a fundamental higher-order function that captures and \emph{reifies} the abstract pattern of structural recursion over lists.

Structural recursion follows a template based on the definition of the data. For a list, the data is either an empty list \texttt{[]} or a value \texttt{x} prepended to another list \texttt{(x:xs)}. The abstract pattern for \texttt{foldr} is defined as:

\begin{itemize}
    \item \textbf{Base Case:}
          \[
              g [] = z
          \]
    \item \textbf{Recursive Step:}
          \[
              g (x:xs) = f\ x\ (g\ xs)
          \]
\end{itemize}

Here, \texttt{z} is the ``zero'' or starting value used for the empty list, and \texttt{f} is a function that combines the current element \texttt{x} with the result of the recursive call on the rest of the list (\texttt{g xs}).

\subsection{How \texttt{foldr} Works}
The name \texttt{foldr} comes from the fact that it associates to the right. When you call

\[
    \texttt{foldr f z [x1, x2, \dots, xn]}
\]

it expands into:

\[
    f\ x_1\ (f\ x_2\ (\dots (f\ x_n\ z) \dots ))
\]

\subsubsection{The ``Replacement'' Perspective}
A useful way to visualize \texttt{foldr} is to think of it as a function that replaces the constructors of a list:

\begin{itemize}
    \item Every \texttt{cons} operator (\texttt{:}) is replaced by the function \texttt{f}.
    \item The empty list (\texttt{[]}) at the end is replaced by the starting value \texttt{z}.
\end{itemize}

For example, the list \texttt{1 : 2 : 3 : []} when processed by \texttt{foldr (+) 0} becomes:

\[
    1 + (2 + (3 + 0))
\]

\subsection{Concrete Examples}
Many standard list functions are simply specific instances of this structural recursion pattern:

\begin{itemize}
    \item \textbf{Summing a list:} \texttt{foldr (+) 0} \\
          Here, \texttt{f} is \texttt{+} and \texttt{z} is \texttt{0}.

    \item \textbf{Calculating length:} \texttt{foldr (\textbackslash x acc -> 1 + acc) 0} \\
          For every element \texttt{x}, it ignores the value and adds \texttt{1} to the accumulated count \texttt{acc}.

    \item \textbf{Mapping:} \texttt{foldr (\textbackslash x acc -> (h x) : acc) []} \\
          This reconstructs the list while applying function \texttt{h} to each element.
\end{itemize}

\subsection{Why Use \texttt{foldr}?}
\begin{itemize}
    \item \textbf{Laziness:} Because \texttt{foldr} is right-associative, it can work with infinite lists if the combining function \texttt{f} is lazy in its second argument. For example:
          \[
              \texttt{foldr (\&\&) True (False : repeat True)}
          \]
          can terminate early because \texttt{False \&\& anything} is \texttt{False}.

    \item \textbf{Reification:} Instead of writing manual recursion for every list-processing function, you ``reify'' the pattern by using \texttt{foldr}, which makes the code more concise and easier to reason about.
\end{itemize}


\section{Left Fold in Haskell (foldl)}

\subsection{The Accumulative Pattern}
In Haskell, \texttt{foldl} (Left Fold) captures the \emph{accumulative recursion} pattern. Unlike standard structural recursion, which processes the head and recurses on the tail, accumulative recursion uses an auxiliary parameter called the \emph{accumulator} to carry partial results through recursive calls.

The abstract pattern for \texttt{foldl} is:

\[
    \begin{aligned}
        \text{Base Case:}      & \quad g\ []\ a = a                    \\
        \text{Recursive Step:} & \quad g\ (x:xs)\ a = g\ xs\ (f\ a\ x)
    \end{aligned}
\]

Here, the function \(f\) combines the current accumulator \(a\) with the current element \(x\) to produce a new accumulator, which is then passed into the recursive call on the tail \(xs\).

\subsection{How \texttt{foldl} Works}
\texttt{foldl} associates to the left. Evaluating

\[
    \texttt{foldl f z [x1, x2, \dots, xn]}
\]

expands as:

\[
    f\ (...(f\ (f\ z\ x_1)\ x_2)\ ...)\ x_n
\]

For example:

\[
    \texttt{foldl (+) 0 [1,2,3]} = ((0 + 1) + 2) + 3
\]

This is fundamentally different from \texttt{foldr}, which expands as:

\[
    1 + (2 + (3 + 0))
\]

\subsection{Comparison: \texttt{foldr} vs. \texttt{foldl}}
While both functions reduce a list to a single value, they differ in associativity, performance, and behavior with infinite lists.

\begin{tabular}{|l|l|l|}
    \hline
    \textbf{Feature}        & \textbf{foldr (Right Fold)}                               & \textbf{foldl (Left Fold)}                                                \\
    \hline
    Associativity           & Right-associative: \(x_1 + (x_2 + (x_3 + z))\)            & Left-associative: \(((z + x_1) + x_2) + x_3\)                             \\
    \hline
    Abstract Pattern        & Structural Recursion: \(g(x:xs) = f\ x\ (g\ xs)\)         & Accumulative Recursion: \(g(x:xs)\ a = g\ xs\ (f\ a\ x)\)                 \\
    \hline
    Infinite Lists          & Can work if \(f\) is lazy in its second argument          & Cannot work; must reach the end before finishing                          \\
    \hline
    Recursion Type          & Builds up expression in memory (not tail-recursive)       & Naturally tail-recursive; may build thunks unless \texttt{foldl'} is used \\
    \hline
    Constructor Replacement & Replaces \texttt{:} with \(f\) and \texttt{[]} with \(z\) & Processes list from seed \(z\) and first element                          \\
    \hline
\end{tabular}

\subsection{Which One to Use?}
\begin{itemize}
    \item Use \texttt{foldr} when working with infinite lists or when producing a new list (e.g., \texttt{map} or \texttt{filter}), because laziness allows early results.
    \item Use \texttt{foldl}, or more commonly its strict version \texttt{foldl'}, for tail-recursive calculations on finite lists, such as summing numbers, to minimize memory usage by updating an accumulator.
\end{itemize}





\section{foldr and foldl expected outputs on the same list}

Gemini said \\
Whether foldr and foldl give the same output depends entirely on the function you are using to fold the list.

The rule of thumb is: If the operation is associative (like addition) and commutative, the results will usually be the same. If the order of operations matters (like subtraction or building a list), the results will be very different.

\subsection{1. When the Output is the Same}
For operations where the order doesn't matter, foldr and foldl are interchangeable in terms of the final result.

\textbf{Addition:} foldr (+) 0 [1,2,3] vs foldl (+) 0 [1,2,3]

\begin{verbatim}
foldr: 1+(2+(3+0))=6
foldl: ((0+1)+2)+3=6
\end{verbatim}

\textbf{Multiplication:} Both will result in 1*2*3=6.

\subsection{2. When the Output is Different}
For operations where the "direction" changes the logic, you will get different results.

\textbf{Example A: Subtraction} \\
Subtraction is not associative, so the grouping changes everything.

\begin{verbatim}
foldr (-) 0 [1, 2, 3]
-- Logic: 1-(2-(3-0))
-- Result: 2

foldl (-) 0 [1, 2, 3]
-- Logic: ((0-1)-2)-3
-- Result: -6
\end{verbatim}

\textbf{Example B: List Construction (The "Reverse" Effect)} \\
foldl naturally flips things because it processes the "head" first and carries it forward, whereas foldr processes the "tail" first.

\begin{verbatim}
-- foldr keeps the order (replacing : with :)
foldr (:) [] [1,2,3] -- Output: [1,2,3]

-- foldl reverses the order (if we flip the arguments)
foldl (\acc x -> x : acc) [] [1,2,3] -- Output: [3,2,1]
\end{verbatim}

\subsection{3. Summary Comparison}

\begin{tabular}{|l|l|l|}
    \hline
    Feature        & foldr (Fold Right)                & foldl (Fold Left)                       \\
    \hline
    Direction      & Right-to-left (starts at the end) & Left-to-right (starts at the beginning) \\
    Association    & Right-associative: f x (f y z)    & Left-associative: f (f z x) y           \\
    Infinite Lists & Can work (if f is lazy)           & Will hang (must reach the end first)    \\
    Memory         & Uses a "call stack"               & Uses an "accumulator"                   \\
    \hline
\end{tabular}

\subsection*{A Note on Performance}
In practice, if you are working with a large list and want the "left" behavior, you should almost always use \texttt{foldl'} (the version with the prime). Standard \texttt{foldl} is often avoided in Haskell because it builds up a massive "thunk" (a pile of unfinished math) in memory, whereas \texttt{foldl'} calculates the result immediately at each step.






\section{Understanding \texttt{lengthList} in Haskell}

Consider the following Haskell code:

\begin{verbatim}
lengthList :: [a] -> Int
lengthList = myFoldr (\_ acc -> 1 + acc) 0
\end{verbatim}

We will break down the key components: the \texttt{\textbackslash} and the \texttt{\_} symbols.

\subsection{The \textbackslash Symbol: Lambda Functions}

In Haskell, \texttt{\textbackslash} introduces an \emph{anonymous function}, also called a \emph{lambda function}. Its general syntax is:

\begin{verbatim}
\arg1 arg2 ... -> expression
\end{verbatim}

In the \texttt{lengthList} function:

\begin{verbatim}
\_ acc -> 1 + acc
\end{verbatim}

is a lambda function equivalent to defining a named function:

\begin{verbatim}
addOneIgnoreX :: a -> Int -> Int
addOneIgnoreX _ acc = 1 + acc
\end{verbatim}

Using a lambda allows the function to be written inline without needing a separate name.

\subsection{The \_ Symbol: Wildcard}

The underscore \texttt{\_} is called the \emph{wildcard} in Haskell. It tells the compiler:

\begin{quote}
    "I do not care about this argument."
\end{quote}

In the lambda \texttt{\_ acc -> 1 + acc}, the first argument represents the current element of the list. We do not need the element itself for counting the length, so we ignore it using \texttt{\_}. This is a common Haskell pattern when only the accumulator matters.

\subsection{How It Works: Step-by-Step}

\texttt{myFoldr} performs a right fold over the list. For each element, it applies the lambda function to the element and the accumulator:

\[
    (\_\, acc \to 1 + acc)\ 5 \ ((\_\, acc \to 1 + acc)\ 7 \ ((\_\, acc \to 1 + acc)\ 9\ 0))
\]

Evaluating from the innermost call:

\begin{align*}
    (\_\, acc \to 1 + acc)\ 9\ 0 & = 1 + 0 = 1 \\
    (\_\, acc \to 1 + acc)\ 7\ 1 & = 1 + 1 = 2 \\
    (\_\, acc \to 1 + acc)\ 5\ 2 & = 1 + 2 = 3
\end{align*}

Thus, \texttt{lengthList [5,7,9] = 3}, which is the length of the list.

\subsection{Summary Table of Symbols}

\begin{tabular}{llp{7cm}}
    \textbf{Symbol}         & \textbf{Concept} & \textbf{Explanation}                                           \\
    \hline
    \texttt{\textbackslash} & Lambda           & Anonymous function, defined inline without a name.             \\
    \texttt{\_}             & Wildcard         & Ignore the argument; used when the value itself is irrelevant. \\
    \texttt{->}             & Lambda body      & Separates the arguments from the expression in the function.   \\
\end{tabular}

\subsection{Key Insight}

Haskell encourages a declarative, concise style. Using \texttt{\_ acc -> 1 + acc} makes it immediately clear that:

\begin{itemize}
    \item The elements of the list are irrelevant for computing length.
    \item Only the accumulator matters.
    \item No extra function names are needed.
\end{itemize}


\section{foldr and foldl expected outputs on the same list}

Gemini said \\
Whether foldr and foldl give the same output depends entirely on the function you are using to fold the list.

The rule of thumb is: If the operation is associative (like addition) and commutative, the results will usually be the same. If the order of operations matters (like subtraction or building a list), the results will be very different.

\subsection{1. When the Output is the Same}
For operations where the order doesn't matter, foldr and foldl are interchangeable in terms of the final result.

\textbf{Addition:} foldr (+) 0 [1,2,3] vs foldl (+) 0 [1,2,3]

\begin{verbatim}
foldr: 1+(2+(3+0))=6
foldl: ((0+1)+2)+3=6
\end{verbatim}

\textbf{Multiplication:} Both will result in 1*2*3=6.

\subsection{2. When the Output is Different}
For operations where the "direction" changes the logic, you will get different results.

\textbf{Example A: Subtraction} \\
Subtraction is not associative, so the grouping changes everything.

\begin{verbatim}
foldr (-) 0 [1, 2, 3]
-- Logic: 1-(2-(3-0))
-- Result: 2

foldl (-) 0 [1, 2, 3]
-- Logic: ((0-1)-2)-3
-- Result: -6
\end{verbatim}

\textbf{Example B: List Construction (The "Reverse" Effect)} \\
foldl naturally flips things because it processes the "head" first and carries it forward, whereas foldr processes the "tail" first.

\begin{verbatim}
-- foldr keeps the order (replacing : with :)
foldr (:) [] [1,2,3] -- Output: [1,2,3]

-- foldl reverses the order (if we flip the arguments)
foldl (\acc x -> x : acc) [] [1,2,3] -- Output: [3,2,1]
\end{verbatim}

\subsection{3. Summary Comparison}

\begin{tabular}{|l|l|l|}
    \hline
    Feature        & foldr (Fold Right)                & foldl (Fold Left)                       \\
    \hline
    Direction      & Right-to-left (starts at the end) & Left-to-right (starts at the beginning) \\
    Association    & Right-associative: f x (f y z)    & Left-associative: f (f z x) y           \\
    Infinite Lists & Can work (if f is lazy)           & Will hang (must reach the end first)    \\
    Memory         & Uses a "call stack"               & Uses an "accumulator"                   \\
    \hline
\end{tabular}

\subsection*{A Note on Performance}
In practice, if you are working with a large list and want the "left" behavior, you should almost always use \texttt{foldl'} (the version with the prime). Standard \texttt{foldl} is often avoided in Haskell because it builds up a massive "thunk" (a pile of unfinished math) in memory, whereas \texttt{foldl'} calculates the result immediately at each step.


\section{Improving a Guess for the Square Root}

Given a number $a$ and a number $g$, where $g$ is a guess for $\sqrt{a}$, we can come up with a new (and improved) guess using the following computation, often referred to as the \textbf{Newton-Raphson method} or the \textbf{Babylonian method} for square roots.

\subsection{Formula}

The new guess $g_{\text{next}}$ is computed as:

\[
    g_{\text{next}} = \frac{g + \frac{a}{g}}{2}
\]

\subsection{Explanation}

This computation works by taking the average of your current guess $g$ and the value obtained by dividing $a$ by that guess $\frac{a}{g}$:

\begin{itemize}
    \item If $g$ is an underestimate, then $\frac{a}{g}$ will be an overestimate.
    \item If $g$ is an overestimate, then $\frac{a}{g}$ will be an underestimate.
\end{itemize}

By taking the average of these two values, the new guess moves significantly closer to the true square root.

In Haskell, this logic is often implemented using the \texttt{iterate} function to generate a sequence of increasingly accurate approximations. For example:

\begin{verbatim}
-- Generates an infinite sequence of improved guesses for sqrt(a)
sqrtApprox a = iterate (\g -> (g + a / g) / 2) initialGuess
\end{verbatim}

This produces an infinite list where each element is a better guess than the last.


% template
\begin{itemize}
    \item X
          \begin{enumerate}
              \item Y
          \end{enumerate}
\end{itemize}

\end{document}